\chapter{Conclusions and Future Work} \label{ch:conclusion}

This chapter presents a short summary of our results, each of which we described in earlier chapters in more detail.

During this thesis, we focused on four main topics and explored each one in-depth.
\begin{itemize}
    \item We examined the feasibility of \textbf{obfuscating the Classic McEliece and HQC KEMs}, and all their parameter sets. If appropriate assumptions from Xagawa~\cite{EC:Xagawa22} hold, then Classic McEliece is a natural OKEM, and an extension of HQC with our proposed filter-encode obfuscator also achieves ciphertext uniformity.
    Even though there are only few practical changes involved at this step, it is important to understand the KEMs we obfuscate and reason about our obfuscators clearly in order to avoid flaws in the OKEM.
    
    \item We then presented an \textbf{implementation of the \drivel{} obfuscated key exchange protocol}, integrated with existing Tor pluggable transports. Some obstacles related to large public key sizes necessitated workarounds. Crucially, we added crypto-agility functionality into the existing repository and thus decoupled the KEM and OKEM implementations from the pluggable transport itself. This thesis specifies 12 parameter sets as predefined combinations of KEMs and OKEMs that minimize the total traffic volume required for a handshake. The flexibility that our implementation now has will help subsequent changes, demonstrating the value of modularization and abstraction.
    
    \item To \textbf{quantify the performance of our implementation}, we used Go benchmarks and observed local deployments with Docker containers. Our evaluation shows that our implementation lacks some optimizations, partly due to the design of the \drivel{} protocol and partly due to our focus on correctness over performance. However, the increase in traffic volume was to be expected due to the use of quantum-safe algorithms instead of traditional ones. The overall traffic pattern remains roughly identical to \obfsfour{}, and bridge operators may make trade-offs regarding latency and memory usage by selecting appropriate parameter sets for \drivel{}. Here we again come back to the idea of providing a flexible protocol that can adapt to future needs, albeit within some sensible preselection of KEM/OKEM combinations.

    \item Finally, we \textbf{suggested changes to the \drivel{} protocol} to enhance its performance and possibly simplify its construction. We believe that the protocol might benefit from leveraging functionality from the framing layer used for the data transfer phase, both for integrity and for more flexible traffic shaping. For these changes, a more in-depth analysis of the framing layer is required, and the implementation would have to undergo some refactoring. This demonstrates the importance of understanding the bigger picture, as some components introduced for an independent purpose might prove useful and reusable.
\end{itemize}

To conclude, we present our views on two main directions for future research.
One possibility for building upon this thesis involves optimizing our implementation of \drivel{}: Its performance is currently suboptimal, and a thorough investigation using Go's profiling tools may reveal various avenues for improvements, in particular for our implementation of ML-Kemeleon. By decreasing the higher performance demands compared to \obfsfour{}, further work could aid the adoption of the protocol. \\
In addition, one could also study the security properties of the framing layer in more detail: Currently, it is unclear what exact guarantees the framing layer should provide in order to make our proposed adaptations of \cref{sssec:variant-framing} possible. By defining these requirements more precisely, a security proof of a simplified \drivel{} variant could be constructed, which uses the framing layer. This would simplify the overall design of the pluggable transport, including both the handshake and the data transfer phase, and might improve the implementation's performance at the same time.
